\section{Introduction}
\label{sec:introduction}

Tabular foundation models (TFMs) learn predictive priors through pretraining
across tasks and perform predictions via in-context learning without
dataset-specific parameter updates
\citep{muller2023tabpfn,hollmann2025tabpfn,qu2025tabicl,ma2025tabdpt}.
Each prediction depends jointly on the query and its labeled context. Changes
in context examples, features, or labels can therefore alter predictions even
when the model and query remain unchanged
\citep{abdullah2026contextsampling}. This dependence allows TFMs to adapt to
new tasks, but also exposes their predictions to unreliable context
information. A natural strategy is to evaluate each query under multiple
contexts, reducing reliance on any single construction. Such contexts can be
obtained through compression, retrieval, sample selection, transformations,
and distribution alignment
\citep{feuer2024tunetables,thomas2024localpfn,xu2025mixturepfn,
treerath2026tabaicl,hollmann2025tabpfn,khalifa2023dense}. These constructions
can provide complementary evidence, but their predictions may disagree and
their reliability varies across queries and conditions. Moreover, when the
prediction from the original context (\emph{Base}) is already accurate,
indiscriminate averaging can make it worse. Improving robustness therefore
requires a query-level decision: whether the alternatives provide sufficient
evidence to depart from Base, and by how much.

We propose \method{} (Adaptive Context Ensembling), a model-agnostic and
query-adaptive inference framework that implements this decision for frozen
TFMs. \method{} evaluates complementary context constructions with the same
predictor, aggregating their predictions via cross-context agreement and
confidence. A task-specific gate then controls the departure from Base,
anchoring the final output to Base when alternative evidence is weak.
Theoretically, we prove a target risk bound guaranteeing that this
evidence-weighted integration effectively limits harmful updates. At runtime,
\method{} requires neither ground-truth target labels nor backbone updates, and
supports reusable context computations to substantially reduce latency. Across
four frozen TFMs on OpenML benchmarks spanning five controlled context
perturbations and seven distribution shifts, as well as synthetic and natural
shifts, \method{} consistently improves prediction robustness while preserving
clean performance. Furthermore, the framework seamlessly transfers to
conventional tabular predictors.

The main contributions of this paper are summarized as follows:
\begin{itemize}
    \item We formulate robust inference for frozen TFMs as deciding when, and
    by how much, alternative context predictions should modify Base.
    \item We introduce \method{}, combining complementary contexts,
    evidence-weighted aggregation, and query-dependent integration with Base,
    and establish sufficient conditions for risk improvement.
    \item We demonstrate robustness under controlled perturbations and
    distribution shifts, supported by synthetic and natural-shift datasets,
    while preserving clean performance.
\end{itemize}
\FloatBarrier
